Лето, письменный стол. Направо дверь. На столе картина. На картине нарисована лошадь, а в зубах у лошади цыган.

Ольга Петровна колет дрова. При каждом ударе с носа Ольги Петровны соскакивает пенсне. Евдоким Осипович сидит в креслах и курит.

ОЛЬГА ПЕТРОВНА (ударяет колуном по полену, которое, однако, нисколько не раскалывается).

ЕВДОКИМ ОСИПОВИЧ. Тюк!

ОЛЬГА ПЕТРОВНА (надевая пенсне, бьёт по полену).

ЕВДОКИМ ОСИПОВИЧ. Тюк!

ОЛЬГА ПЕТРОВНА (надевая пенсне). Евдоким Осипович! Я вас прошу, не говорите этого слова <<тюк>>.

ЕВДОКИМ ОСИПОВИЧ. Хорошо, хорошо.

ОЛЬГА ПЕТРОВНА (ударяет колуном по полену).

ЕВДОКИМ ОСИПОВИЧ. Тюк!

ОЛЬГА ПЕТРОВНА (надевая пенсне). Евдоким Осипович! Вы обещали не говорить этого слова <<тюк>>.

ЕВДОКИМ ОСИПОВИЧ. Хорошо, хорошо, Ольга Петровна! Больше не буду.

ОЛЬГА ПЕТРОВНА (ударяет колуном по полену).

ЕВДОКИМ ОСИПОВИЧ. Тюк!

ОЛЬГА ПЕТРОВНА (надевая пенсне). Это безобразие! Взрослый пожилой человек и не понимает простой человеческой просьбы!

ЕВДОКИМ ОСИПОВИЧ. Ольга Петровна! Вы можете спокойно продолжать вашу работу. Я больше мешать не буду.

ОЛЬГА ПЕТРОВНА. Ну я прошу вас, я очень прошу вас: дайте мне расколоть хотя бы это полено.

ЕВДОКИМ ОСИПОВИЧ. Колите, конечно, колите!

ОЛЬГА ПЕТРОВНА (ударяет колуном по полену).

ЕВДОКИМ ОСИПОВИЧ. Тюк!

Ольга Петровна роняет колун, открывает рот, но ничего не может сказать.

Евдоким Осипович встаёт с кресел, оглядывает Ольгу Петровну с головы до ног и медленно уходит.

Ольга Петровна стоит неподвижно с открытым ртом и смотрит на удаляющегося Евдокима Осиповича.

\begin{center}
Занавес медленно опускается.
\end{center}

\begin{flushright}
<\dots{}>
\end{flushright}

\begin{center}
* * *   
\end{center}
